\documentclass[11pt]{article}
\usepackage{ap-psych}
\title{Unit 4: Social Psychology and Personality: Question Set}
\author{Paul Wang}
\date{November 2026}
\docpart{Unit 4 \textperiodcentered\ Question Set}

\begin{document}
\maketitle

\begin{abstract}
Exam-style multiple choice on all seven Unit 4 topics, with a shared-study set on conformity and a full Article Analysis Question on group size and helping online. Nearly every item is a scenario; the distractors are the terms most often confused with the right answer (fundamental attribution error vs.\ self-serving bias, group polarization vs.\ groupthink, James--Lange vs.\ two-factor). All items and the AAQ source are original and fictional, written for practice. Allow about 30 minutes for the multiple choice and 25 minutes for the AAQ.
\end{abstract}

\scoringline{MCQ: \blank\ / \qtotal \qquad AAQ: \blank\ / 7 \qquad Target: 21+ MCQ and 5+ AAQ \qquad Pace: 72 s per MCQ}

\section{Social thinking (4.1--4.2)}

\begin{q}{1}{4.1 Actor-observer}{Own behavior \(\rightarrow\) situation; others' \(\rightarrow\) disposition.}
Pia says she snapped at a coworker because she was exhausted, but says her coworker snaps because he is rude. This is the
\chh{actor--observer bias}{just-world phenomenon}{mere-exposure effect}{false consensus effect}
\end{q}

\begin{q}{3}{4.1 Explanatory style}{Internal, stable, global explanation of a setback = pessimistic style.}
After one failed quiz, Leo says, ``I'm stupid and I'll fail at everything.'' His explanation is
\ch{external, temporary, and specific}{an internal locus of control applied optimistically}{internal, stable, and global: a pessimistic explanatory style}{an example of downward social comparison}
\end{q}

\begin{q}{4}{4.1 Mere exposure}{Repeated exposure increases liking.}
A song that Maya disliked at first grows on her after hearing it daily on the radio. This is the
\chh{halo effect}{overjustification effect}{self-fulfilling prophecy}{mere-exposure effect}
\end{q}

\begin{q}{2}{4.1 Relative deprivation}{Feeling worse off compared with a reference group despite gains.}
After a raise, Sam feels resentful because his friends at another firm earn more. He is experiencing
\chh{the self-serving bias}{relative deprivation}{reaction formation}{deindividuation}
\end{q}

\begin{q}{3}{4.2 Prejudice}{Stereotype = belief, prejudice = attitude, discrimination = behavior.}
A landlord refuses to rent to applicants from a particular ethnic group. This behavior is best labeled
\chh{a stereotype}{prejudice}{discrimination}{out-group homogeneity}
\end{q}

\begin{q}{1}{4.2 Belief perseverance}{Keeping a belief after its basis is discredited.}
Even after learning that the study linking vaccines to autism was fraudulent and retracted, some parents continue to believe the link. This is
\chh{belief perseverance}{cognitive dissonance}{the halo effect}{informational influence}
\end{q}

\begin{q}{2}{4.2 Foot-in-the-door}{Small request first, then a larger one = foot-in-the-door.}
A campaign worker first asks residents to sign a petition; a week later, those who signed are more likely to agree to put a large sign on their lawn. This is the
\ch{door-in-the-face technique}{foot-in-the-door technique}{peripheral route}{reciprocity norm}
\end{q}

\begin{q}{4}{4.2 Out-group homogeneity}{Seeing ``them'' as all alike = out-group homogeneity.}
Fans of one school's team say, ``Students at our school are all different, but the kids at our rival school are all the same.'' This is
\chh{ethnocentrism}{group polarization}{the just-world phenomenon}{out-group homogeneity}
\end{q}

\section{Social situations (4.3)}

\begin{stimulus}
\textbf{Study E.} In a line-judgment task like Asch's, 80 college students answered aloud after confederates who gave unanimous wrong answers. In a second condition, one confederate gave the correct answer. Percent of trials on which participants conformed:
\begin{center}
\begin{tabular}{lc}
\toprule
Condition & Conformity\\\midrule
Unanimous wrong majority & 34\%\\
One dissenter present & 9\%\\
Unanimous majority, answers written privately & 12\%\\
\bottomrule
\end{tabular}
\end{center}
\end{stimulus}

\begin{q}{3}{4.3 Conformity}{A single dissenter breaks unanimity and sharply reduces conformity.}
Which conclusion is best supported by Study E?
\ch{Conformity is unaffected by group unanimity.}{Private answers increase conformity.}{A single ally who breaks the group's unanimity sharply reduces conformity.}{Most participants conformed on most trials.}
\end{q}

\begin{q}{2}{4.3 Normative influence}{Conformity drops when answers are private: the pressure was about acceptance (normative).}
The drop in conformity when answers were written privately suggests that conformity in the public condition was driven mainly by
\ch{informational social influence}{normative social influence}{social facilitation}{groupthink}
\end{q}

\begin{q}{1}{4.3 Obedience}{Following a direct order from an authority = obedience (Milgram).}
In Milgram's studies, which change most reduced obedience?
\ch{Two peers who refused to continue}{Making the experimenter wear a lab coat}{Holding the study at Yale}{Placing the learner in a separate room}
\end{q}

\begin{q}{4}{4.3 Groupthink}{Harmony-seeking suppresses dissent and leads to a poor decision.}
A product team, eager to avoid conflict with its confident leader, approves a flawed launch plan without voicing doubts. This is
\chh{group polarization}{social loafing}{deindividuation}{groupthink}
\end{q}

\begin{q}{3}{4.3 Deindividuation}{Anonymity in a crowd lowers self-awareness and restraint.}
Masked, anonymous fans in a large crowd vandalize property they would never damage alone. This is best explained by
\chh{social facilitation}{the reciprocity norm}{deindividuation}{altruism}
\end{q}

\begin{q}{1}{4.3 Superordinate goals}{A shared goal requiring cooperation reduces intergroup conflict.}
Two rival clubs stop feuding after they must work together to raise money to save the school's only bus. This illustrates
\chh{a superordinate goal}{a social trap}{group polarization}{the false consensus effect}
\end{q}

\section{Personality (4.4--4.5)}

\begin{q}{2}{4.4 Defense mechanisms}{Returning to childish behavior under stress = regression.}
After his baby sister is born, a 6-year-old begins sucking his thumb and wetting the bed again. Freud would call this
\chh{sublimation}{regression}{projection}{displacement}
\end{q}

\begin{q}{4}{4.4 Psychodynamic}{Ego mediates between the id's demands and the superego's morality.}
In Freud's theory, the part of personality that seeks realistic ways to satisfy desires while respecting moral limits is the
\chh{id}{superego}{unconscious}{ego}
\end{q}

\begin{q}{3}{4.4 Projective tests}{Ambiguous pictures to tell stories about = TAT.}
A psychologist shows a client ambiguous pictures of people and asks her to make up a story about each. This is the
\ch{MMPI}{Big Five inventory}{Thematic Apperception Test}{Implicit Association Test}
\end{q}

\begin{q}{1}{4.5 Trait assessment}{Self-report inventories are more reliable and valid than projective tests.}
Compared with projective tests, personality inventories such as Big Five questionnaires are generally
\ch{more reliable and valid}{better at revealing unconscious conflicts}{based on psychodynamic theory}{scored by the clinician's interpretation of stories}
\end{q}

\begin{q}{2}{4.5 Big Five}{Anxious, moody, easily upset = high neuroticism.}
A person who worries constantly, is easily upset, and has frequent mood swings would score high in
\chh{agreeableness}{neuroticism}{openness}{conscientiousness}
\end{q}

\section{Motivation and emotion (4.6--4.7)}

\begin{q}{4}{4.6 Drive reduction}{A need creates a drive that pushes behavior to restore homeostasis.}
After a long run, Zoe feels thirsty and drinks until the feeling goes away. This best fits
\chh{arousal theory}{incentive theory}{instinct theory}{drive-reduction theory}
\end{q}

\begin{q}{1}{4.6 Self-determination}{Autonomy, competence, relatedness = self-determination theory.}
A coach lets players choose drills, gives feedback that builds skill, and creates a close team atmosphere. She is applying
\ch{self-determination theory}{drive-reduction theory}{the Yerkes--Dodson law}{instinct theory}
\end{q}

\begin{q}{3}{4.6 Conflicts}{Choosing between two unpleasant options = avoidance-avoidance.}
Ty must either clean the garage or mow the lawn, and he dislikes both. He faces
\ch{an approach--approach conflict}{an approach--avoidance conflict}{an avoidance--avoidance conflict}{a multiple approach--avoidance conflict}
\end{q}

\begin{q}{2}{4.7 Cannon-Bard}{Arousal and emotion occur at the same time and independently.}
A theory holds that the heart pounding and the feeling of fear occur at the same moment, each triggered by the thalamus. This is
\ch{the James--Lange theory}{the Cannon--Bard theory}{the two-factor theory}{the facial-feedback hypothesis}
\end{q}

\begin{q}{4}{4.7 Facial feedback}{Making a facial expression intensifies the matching emotion.}
Participants who held a pen between their teeth (forcing a smile) rated cartoons as funnier. This supports
\ch{the James--Lange theory only}{display rules}{broaden-and-build theory}{the facial-feedback hypothesis}
\end{q}

\section{Article Analysis Question}

\begin{frq}{Question 1 \textperiodcentered\ AAQ \textperiodcentered\ 7 points \textperiodcentered\ 25 minutes}
\small
\textbf{Source (summarized; fictional, for practice).} Delacroix, M., \& Iqbal, S. (2024). Audience size and helping in online communities. \emph{Journal of Social Behavior Online}, 12, 88--104.

\textbf{Participants.} 300 adult volunteers (ages 18--40) recruited through a university study pool in France. All signed consent forms describing the study as ``research on online teamwork.'' Because the true purpose was concealed, every participant was debriefed afterward and offered the chance to withdraw their data.

\textbf{Method.} Participants were randomly assigned to join a chat room that displayed 1, 5, or 20 other ``members'' (all simulated). During a team task, a simulated member typed that her laptop was overheating and asked for someone to tell the moderator. Helping was defined as sending a message to the moderator within 2 minutes.

\textbf{Results.} Percent who helped: 1 other member, 71\%; 5 others, 38\%; 20 others, 14\%. Group size and helping correlated \(r = -.49\).

\medskip
\begin{enumerate}[label=(\Alph*),leftmargin=2em,itemsep=3pt]
\item Identify the research method used in the study.
\item State the operational definition of helping.
\item Describe what the correlation of \(r = -.49\) indicates.
\item Identify at least one ethical guideline applied by the researchers.
\item Explain the extent to which the findings may or may not be generalizable, using specific information from the study.
\item Explain how at least one finding supports or refutes the claim that ``people are less likely to help when they believe many others could help.''
\end{enumerate}
\end{frq}

\makekey

\subsection*{AAQ scoring guide (7 points)}
\small
\begin{compare}{Part}{Earns the point}{Does not earn the point}
A (1) & ``Experiment'' (group size was manipulated with random assignment). & ``Correlational,'' even though a correlation is reported. \\
B (1) & Sending a message to the moderator within 2 minutes of the request. & ``Being helpful''; ``caring about the member.'' \\
C (1) & A moderate negative relationship: as the number of other members increased, the percent who helped decreased. & ``Strong positive''; ``49\% helped.'' \\
D (1) & Debriefing after deception, informed consent, \emph{or} right to withdraw data, linked to the study. & ``They followed IRB rules'' with no specifics. \\
E (1) & Limited: 18--40-year-old French university-pool volunteers in a simulated chat may not reflect older adults, other cultures, or real emergencies. \emph{Or} supported: large sample and random assignment support generalizing to similar online settings. & ``Only 300 people'' with no reasoning. \\
F (2) & 1 pt: supports the claim, citing a result (71\% with 1 other vs.\ 14\% with 20 others). 1 pt: explains with a concept, e.g., \emph{diffusion of responsibility} / the \emph{bystander effect}: as more potential helpers are present, each person feels less personally responsible. & Naming the bystander effect without the numbers, or citing numbers with no concept. \\
\end{compare}

\end{document}
